\begin{abstract}
A compositional representation lets a system understand a new combination by reading its familiar parts back out of the whole. Machine learning now calls almost any successful recombination of familiar material compositional, and this use cannot separate a representation that keeps its parts from one that blends them. We propose that a representation is compositional when a fixed procedure of a declared kind, such as an affine map, reads every part back in its role, including parts on which the procedure was never fitted, and when each part enters every role through one shared code. We prove that a lookup decodes almost any combination of a finite vocabulary, that the rank of a linear relation and the width of the state bound the recovery of fresh parts, and that a shared code fixes every role reversal from a single composite. Analytic operators meet these bounds. In networks trained on glyph composites and in a learned recursive composer, the states keep fresh parts recoverable, while the networks' own decoders fail where training did not reach. We then measure the same quantities in 28 pretrained models, from text encoders to language models with 15 billion parameters and image encoders. Training lifts every model above a blend, yet no fully trained model gets more than 60\% of the way to a role--filler reference. Recovery in Pythia-1.4b peaks early in training and then declines, and every language model binds roles to word order. A system's states and the procedures that read them therefore become compositional separately, and current scale and training leave both incomplete.
\end{abstract}

%=====================================================================
\section{Introduction}
\label{sec:intro}

Composition is central to human intelligence. People understand sentences they have never heard by building their meaning from familiar parts \citep{frege1923}, and \citet{fodor1988} argued that representations with constituent structure explain this ability. Deep learning studies the same ability as compositional generalization, the production of correct outputs for novel combinations of familiar atoms \citep{lake2018,keysers2020,kim2020cogs}. Text-to-image models compose concepts into new scenes \citep{ramesh2021,okawa2023}, yet they often misplace objects \citep{conwell2022} and attach an attribute to the wrong object \citep{rassin2022}. The key question is therefore which representations deserve the name compositional.

Work on this question often starts from a principle about meaning. The meaning of a complex is a function of the meanings of its parts and the way they combine \citep{montague1970}. Every system whose state depends on its inputs meets this principle, so the principle treats an avocado on a chair and an avocado-shaped armchair alike (Figure~\ref{fig:ops}, left). Recent accounts measure the simplicity of a symbolic redescription \citep{elmoznino2024}, the structure a whole adds beyond its parts \citep{elmoznino2026}, or the fit of an explicit composition of inferred primitives \citep{andreas2019}. Identifiability results give sufficient \citep{lachapelle2023,wiedemer2024} and necessary \citep{uselis2026} conditions for compositional generalization. Despite this progress, these accounts leave open when recovering parts from a state shows structure, and when it shows only a familiar vocabulary.

Reading parts back offers a complementary perspective. \citet{vangelder1990} asked for general procedures that decompose a complex into its constituents, and \citet{fodor2002} argued that understanding a complex includes understanding its parts. However, a readout fitted on a few familiar parts decodes almost any combination of them, blends included. This leaves a gap between the parts a system has seen and the parts its states can carry. A fuller account should therefore test recovery on parts that the reading procedure has never seen.

Large pretrained models are a natural place for this test. Probing studies read their states routinely, nouns and photographs supply hundreds of parts, and random weights and training checkpoints supply controls. We can therefore ask whether scale and training make a representation compositional, and whether the roles a model encodes follow meaning or word order. We test the definition on analytic operators, networks trained from scratch, a recursive composer and 28 pretrained models, and then ask whether trained networks use the parts their states carry.

\textbf{Contributions.} We define a representation as compositional when a declared readout recovers every part in its role from composites of parts fresh to the readout, and when each part enters every role through one shared code. We prove that a lookup decodes almost any relation on a finite vocabulary, that the rank of a linear relation and the width of a state bound the recovery of fresh parts, and that a shared code fixes every role reversal from one composite. We show in networks trained from scratch that the states keep fresh parts while the networks' own decoders fail where training did not reach. We show in text encoders, language models with up to 15 billion parameters and image encoders that training lifts every model above a blend, that no fully trained model gets past 60\% of the way to a role--filler reference, and that language models bind roles to word order.

%=====================================================================
\section{Related work}
\label{sec:related}

\textbf{Identifiability.} Additive decoders, together with an encoder that inverts them on recombined slots, guarantee that object-centric models generalize to unseen combinations \citep{brady2023,lachapelle2023,wiedemer2024}. Our role-factorization is their additivity in affine coordinates, and our recoverability is a readout-side, affine analogue of their consistency condition. \citet{uselis2026} argue that compositional generalization in vision embeddings requires linear, orthogonal factors. We use such conditions to define a property of states, ask only for a full-rank role map, and separate recovery by a readout fitted on many parts from recovery by procedures trained on a system's own vocabulary.

\textbf{Probing and decodability.} Neuroscience tests whether a decoder transfers to held-out conditions \citep{bernardi2020}, and linear decodability of constituents accompanies compositional generalization as models scale \citep{redhardt2025}. Decodable factors do not ensure that a decoder generalizes \citep{montero2021,schott2022}, and agents can generalize without compositional messages \citep{chaabouni2020}. Discovered tensor-product decompositions \citep{mccoy2019,soulos2020}, binding identifiers \citep{feng2024} and interchange interventions \citep{geiger2021,geiger2024das} can supply the part codes our definition needs. We add the requirement that a readout transfer to parts outside the vocabulary it was fitted on.

\textbf{Compression.} \citet{elmoznino2024} measure the description length of a representation under a symbolic redescription, and \citet{elmoznino2026} measure the structure a whole adds beyond its parts. Neither account asks that parts return from the whole. A superposition over a small vocabulary compresses well and passes clean-up decoding, and our criterion separates it on fresh parts.

\textbf{The classical debate.} \citet{fodor1988} argued that systematicity requires constituent structure, and \citet{fodor1990} rejected the view of \citet{smolensky1990} that recoverable tensor-product constituents suffice, because such constituents need not act in processing. Our criterion sides with Smolensky on the content of a representation and turns the objection into a test, an intervention on one role's component. The demand that the architecture itself guarantee systematicity lies outside any criterion on states \citep{russin2024,goodale2026}.

%=====================================================================
\section{Recoverable composition}
\label{sec:definition}

In this section, we define recoverable and role-factorized relations and state the results that follow from them. Appendix~\ref{app:proofs} gives the proofs.

\begin{figure}[tbp]
  \centering
  \begin{minipage}[c]{0.5\linewidth}\includegraphics[width=\linewidth]{fig1_avocado.pdf}\end{minipage}\hfill
  \begin{minipage}[c]{0.47\linewidth}\includegraphics[width=\linewidth]{fig_fibers.pdf}\end{minipage}
  \caption{\textbf{Blending and role--filler composition keep different amounts of their parts.} \textbf{Left.} Role--filler composition places the avocado in a role of \emph{on}, so a reader can recover either part. Blending merges the two into an avocado-shaped armchair \citep{openai2021dalle}, from which neither part returns; Canva's text-to-image tool generated the images. \textbf{a.} Under the sum $x + y$, every pair on a line gives the same composite. The vocabulary pairs on one line (dark points) reverse each other's roles, so a lookup over familiar composites finds the summed parts, while fresh pairs (diamonds) can lie anywhere on the line. \textbf{b.} A role--filler relation gives each composite one decomposition.}
  \label{fig:ops}
\end{figure}

\subsection{Parts, composites and readouts}
\label{sec:intuition}

A part is a draw from a distribution $\D$ over possible parts, such as glyphs, nouns or photographs. A system assigns each part $x$ a \emph{part code} $e(x) \in \R^d$, and it assigns a composite $r(x_1, \ldots, x_k)$ of $k$ parts in the roles of a relation $r$ a state $h = g_r(x_1, \ldots, x_k) \in \R^m$. Training exposes a system to a finite \emph{vocabulary} $\Voc$ of parts. A part is \emph{fresh} to a procedure when no composite containing it helped fit the procedure. A glyph can be fresh both to a readout and to the network that encodes it, while a noun in a pretrained model can be fresh only to the readout. A \emph{readout class} $\U$, such as affine maps, contains the procedures allowed to recover a part's code from a state.

Suppose a system adds two part vectors $x$ and $y$. Every pair on the line $x + y = z$ yields the same composite (Figure~\ref{fig:ops}a), so no procedure can return the parts. A tensor-product binding $x \otimes \rho_1 + y \otimes \rho_2$ with independent role vectors keeps one decomposition for each composite (Figure~\ref{fig:ops}b). A small vocabulary hides the difference, because few vocabulary pairs share a line and a lookup finds the summed parts. A fresh pair exposes it, because its line holds infinitely many other decompositions. Two lemmas make this general (Appendix~\ref{app:proofs}). When the coordinates of an additive decomposition are unrestricted, every invertible relation is additive (Lemma~\ref{lem:psi}). On a finite vocabulary, a lookup table inverts every relation that is injective on the vocabulary (Lemma~\ref{lem:lookup}). A definition must therefore fix the readouts in advance and score them on fresh parts.

\subsection{The definition}

\begin{defbox}{Recoverable and role-factorized relations}{recoverable}
Fix a part distribution $\D$, a part code $e$, and a readout class $\U$. A relation $r$, which a system realizes by $g_r$, is \emph{$\varepsilon$-recoverable} if for every role $i$ some $u_i \in \U$ satisfies, for parts $x_1, \ldots, x_k$ drawn independently from $\D$,
\begin{equation}
  \frac{\mathbb{E}\,\lVert u_i(g_r(x_1,\ldots,x_k)) - e(x_i) \rVert^2}{\mathbb{E}\,\lVert e(x) - \mathbb{E}\,e(x) \rVert^2} \;\le\; \varepsilon .
  \label{eq:recover}
\end{equation}
The relation is \emph{$\delta$-role-factorized} if, for some code $e'$ of dimension $d'$ of which $e$ is an affine function, some $c \in \R^m$ and injective matrices $B_1, \ldots, B_k \in \R^{m \times d'}$, one per role and shared by all parts, satisfy
\begin{equation}
  g_r(x_1, \ldots, x_k) = c + \textstyle\sum_{i} B_i\, e'(x_i) + \eta(x_1, \ldots, x_k), \qquad \mathbb{E}\,\lVert \eta \rVert^2 \le \delta\, \mathbb{E}\,\lVert g_r - \mathbb{E}\,g_r \rVert^2 .
  \label{eq:factor}
\end{equation}
We estimate both population quantities by fitting $u_i$, $c$ and $B_i$ on composites of a \emph{probe pool} of parts and evaluating them on composites of other parts, with errors in a declared inner product on part codes.
\end{defbox}

Recoverability asks a fixed procedure of a declared kind to return each part in its role across the whole part distribution, and $1 - \varepsilon$ is the population $R^2$ of the best readout in $\U$. A readout fitted on a small vocabulary is undetermined outside the codes it saw (Proposition~\ref{prop:span}), and scoring it on that vocabulary inflates it (Lemma~\ref{lem:lookup}). Fitting on a large probe pool and scoring on fresh parts removes both biases. Role-factorization asks each part to contribute the same code to every whole, transformed only by its role, the classical demand that a constituent contribute the same thing to every complex containing it \citep{fodor1988}. The maps $B_i$ must be injective, because otherwise $e'$ could hold a separate code for each role and two unrelated lookup tables would pass. A relation fails recoverability when it discards information about its parts, and it fails role-factorization when a part's contribution depends on its co-parts, as \emph{red} does in \emph{red hair} and \emph{red apple} \citep{kamp1995}.

\begin{defbox}{Representational compositionality}{repcomp}
A system is \emph{representationally compositional} with respect to a relation $r$, relative to a declared readout class $\U$ and a declared class $\M$ of methods for finding part codes, to the degree that, for some part code $e$ found under $\M$ that is informative about the parts (Definition~\ref{def:informative}), its realization of $r$ is $\varepsilon$-recoverable and $\delta$-role-factorized. We report the degree by $\delta$ and by
\begin{equation}
  \kappa = \frac{\varepsilon_{\mathrm{blend}} - \varepsilon}{\varepsilon_{\mathrm{blend}} - \varepsilon_{\mathrm{ref}}},
  \label{eq:kappa}
\end{equation}
where the same procedure gives $\varepsilon_{\mathrm{ref}}$ for the concatenation and $\varepsilon_{\mathrm{blend}}$ for the sum of the same part codes. The system is representationally compositional at tolerance $\tau$ if, for some such $e$ informative at $\tau$, $\kappa \ge 1 - \tau$ and $\delta \le \tau$.
\end{defbox}

\begin{defbox}{Informative part codes and productivity}{informative}
A part code $e$ is \emph{informative} at tolerance $\tau$ if an affine map of $e$ recovers a declared description $z(x)$ of the parts that does not depend on the system, with $R^2 \ge 1 - \tau$ on fresh parts. A relation is \emph{productive to depth $n$} if Definition~\ref{def:repcomp} holds when the parts are themselves composites of depth up to $n - 1$, with the state of each composite as its part code.
\end{defbox}

We call a system whose states meet Definition~\ref{def:repcomp} representationally compositional, and a system whose own procedures also read its parts from composites of fresh parts a \emph{compositional system} (Table~\ref{tab:twobytwo}). For a symmetric relation of identically distributed parts, no readout can tell the roles apart, so in population $\varepsilon_{\mathrm{blend}} = 1/2$, $\varepsilon_{\mathrm{ref}} = 0$ and $\kappa = 1 - 2\varepsilon$. The references correct the estimates for finite samples and readout fitting. The definition depends on $\U$ and $\M$ because a decomposition must be found by search, and a negative result reports that the declared methods found none. Unrestricted readouts let any injective system pass, as unrestricted probes \citep{hewitt2019} and causal alignment maps \citep{sutter2025} do. The informativeness clause closes a second loophole. The blend $(Px, Py)$, with $P$ a projection of rank $r < d$, discards $d - r$ dimensions of each part, yet with the code $e(x) = Px$ it is exactly recoverable and role-factorized. We report informativeness wherever a description of the parts exists apart from the system, which holds for glyph latents and fails for nouns and photographs.

\begin{table}[tbp]
  \caption{\textbf{Behavioral and representational compositionality vary independently} \citep[cf.][]{vegner2025}. Rows say whether a declared readout recovers parts from composites of fresh parts, and columns say whether the system's own behavior succeeds on new combinations and new parts.}
  \label{tab:twobytwo}
  \centering\small
  \begin{tabularx}{\linewidth}{@{}P{3.2cm}YY@{}}
    \toprule
    & Behavior succeeds & Behavior fails \\
    \midrule
    Parts recoverable from fresh composites & Compositional system; the relation fixes reversals and new role assignments & Failure of performance; the state holds the parts and the decoder misses them (cf.\ \citealt{montero2021,park2024hidden}) \\
    Parts recoverable only from familiar composites & Generalization beyond the declared readouts; the readout class may be too weak & Neither; the parts and the behavior both fail \\
    \bottomrule
  \end{tabularx}
\end{table}

\subsection{Four consequences}
\label{sec:consequences}

Appendix~\ref{app:proofs} proves two lemmas and ten propositions, and four results carry the argument.

\textbf{Blends lose a computable share of their parts.} A linear relation $g(x_1, \ldots, x_k) = A x$ with $A = [A_1 \cdots A_k]$ gives each composite an affine set of decompositions of dimension $kd - \mathrm{rank}\,A$ (Proposition~\ref{prop:fiber}). For Gaussian parts, no readout of any kind recovers role $i$ with $R^2$ above $\mathrm{tr}(A_i^{\top}(AA^{\top})^{+}A_i)/d$, whose mean over roles is $\mathrm{rank}\,A/(kd)$ (Proposition~\ref{prop:ceiling}). The ceiling is 0.8 for role~1 of $x + 0.5y$, 0.5 for $x + y$, and $m/(kd)$ for a generic projection into $m < kd$ dimensions. For any relation and part distribution, an affine readout of an $m$-dimensional state explains at most the share of the parts' variance in the top $m$ principal components of the stacked part codes (Proposition~\ref{prop:ey}).

\textbf{Familiar parts decode anyway.} A non-invertible relation can be injective on a finite vocabulary, and clean-up then decodes its parts \citep{plate2003,frady2018}. Any procedure learned from a vocabulary shares this limit, whether an experimenter's readout or a network's own decoder. An affine procedure fitted on composites is determined only on the codes it saw, so it covers fresh composites only when the codes of the training parts in each role have full rank $d$ (Proposition~\ref{prop:span}). This result is a readout-side form of the argument of \citet{marcus1998} that networks fail to extend a relation outside their training space.

\textbf{A shared code fixes reversals from one composite.} If a relation is role-factorized with $\delta = 0$ and $[B_1 \cdots B_k]$ has full column rank, an affine function of a single composite gives every role reversal and every assignment of a known part to a new role (Proposition~\ref{prop:systematic}). This is the unbinding argument for tensor-product representations \citep{smolensky1990}. \citet{fodor1990} objected that recoverable constituents need not act in processing. An interchange intervention that replaces one role's component of a state tests this objection, because it should change the system's behavior for that role alone. Without role-factorization, observed composites constrain unobserved ones only when the observed pairs connect every part to every role (Proposition~\ref{prop:separable}).

\textbf{Productivity needs invertibility on the relation's own outputs.} A nested linear blend of depth $n$ into a state of width $d$ loses at least $(2^n - 1)d$ dimensions of its leaves (Proposition~\ref{prop:depth}). A relation that inverts on composites of parts can still fail on composites of composites, because they come from another distribution.

%=====================================================================
\section{Experimental design}
\label{sec:experiments}

In this section, we describe the protocol every experiment shares and the systems we test (Figure~\ref{fig:protocol}).

\begin{figure}[tbp]
  \centering
  \includegraphics[width=\linewidth]{fig_protocol.pdf}
  \caption{\textbf{One protocol measures recoverability in every system.} \textbf{a.} We split the parts, glyphs, nouns or photograph classes, before fitting anything. A probe pool fits readouts and role maps, validation parts choose penalties, layers and part codes, and fresh parts give every reported number. \textbf{b.} Affine readouts fitted on composites of the pool return each role's part code from the system's state, and their error on fresh parts is $\varepsilon$. The same pipeline, applied to the concatenation and to the sum of the same codes, calibrates $\varepsilon$ into the score $\kappa$, which is 1 for role--filler composition and 0 for a blend. \textbf{c.} Role maps give the residual $\delta$ and predict role reversals, and interchange interventions pass edited states through the system's own decoder. Readouts fitted on the pool and on the system's training parts meet the system's behavior on fresh parts. \textbf{d.} A converse paraphrase keeps the figure (the tulip) and reverses the order of mention. A readout of the figure fitted on one phrasing returns the tulip on the other if roles follow meaning, and the barrel if they follow word order. \textbf{e.} The systems we test and their controls.}
  \label{fig:protocol}
\end{figure}

\paragraph{Protocol.} We split the parts before fitting anything. Readouts and role maps fit composites of a probe pool, validation parts choose ridge penalties, layers and part codes, and fresh parts give every reported number. A ridge readout of the system's state returns each role's part code, and its error $\varepsilon$ on fresh parts measures recovery. The identical pipeline, applied to the concatenation and to the sum of the same part codes, calibrates $\varepsilon$ into $\kappa$ (Eq.~\ref{eq:kappa}). Role maps fitted on the pool give the residual $\delta$. For systems that act, we compare a \emph{population readout} fitted on the pool with a \emph{vocabulary readout} fitted on the system's own training parts, and both with the system's behavior on fresh parts. Every number is a mean over seeds, and Appendix~\ref{app:protocol} gives full details.

\paragraph{Analytic operators.} Parts are Gaussian vectors in $\R^{16}$, and fourteen operators combine them in pairs. They include role--filler binding (concatenation and a tensor product), fixed-width superposition (random projections into $0.5$ to $1.5$ times $kd = 32$ dimensions, HRR and MAP binding, and $x + 0.5y$), symmetric blends ($x + y$ and $x \odot y$), and a random $\tanh$ layer. We fit ridge, least-squares and MLP readouts of the role-1 part on composites of a vocabulary of $V$ parts, with 5 seeds for each $V$.

\paragraph{Glyph networks.} A fixed random generator renders a latent $z \in \R^{8}$ as an $8 \times 8$ glyph, and a composite places two glyphs on an $8 \times 16$ canvas under one of three relations. \textsc{beside} uses disjoint regions, \textsc{over} draws both glyphs in one region with the second mirrored, and \textsc{occlude} draws the second at half scale over the first, a non-additive relation. An encoder writes a state $h \in \R^m$, and a decoder renders the same parts in the same roles under a target relation. The state is therefore the system's representation, and the decoder is its own reading procedure. We train an MLP, a CNN, a transformer and a role--filler decoder that decodes two halves of $h$ with one shared glyph decoder, at widths $m$ from 8 to 64 and vocabularies $V$ from 8 to 64 glyphs, with 9 seeds. A quarter of the vocabulary never occupies role~2 in training. For the structural tests, we declare the top 16 principal components of glyph pixels as the part code and retrain each network with 5 seeds.

\paragraph{Recursive composers.} Balanced binary trees of depth $n$ have Gaussian leaves in $\R^8$. A recursive auto-associative memory \citep{pollack1990} applies one learned composer at every node and one learned decomposer recursively, and we train both on depths 1 and 2. Holographic reduced representations (HRR) of the same width $w$ use one pair of unitary role vectors per level. A population readout maps the root state to all $2^n$ leaves, with 5 seeds for each $w$ from 16 to 256.

\paragraph{Text encoders.} Sentences place two nouns in the roles of five asymmetric relations, such as ``a $x$ on a $y$'', and of the control ``a $x$ and a $y$''. We split 305 concrete nouns into 165 for fitting readouts, 40 for validation and 100 fresh to the readout. We test eleven encoders, three from CLIP \citep{radford2021clip}, including those that condition Stable Diffusion 1.5 and 2 \citep{rombach2022}, three from T5 \citep{raffel2020t5}, two from BERT \citep{devlin2019bert} and three from GPT-2 \citep{radford2019gpt2}. A part code is the state of ``a $x$'', the state of ``there is a $x$ here'', or the input embedding, which we reduce to 16 principal components and whiten. Whitening matters because a few directions dominate transformer states \citep{ethayarajh2019,timkey2021}, and the first component holds 99\% of the variance of GPT-2's early-layer code. Each architecture with random weights serves as a control.

\paragraph{Language models.} We apply the same sentences and pipeline to Qwen3-8B \citep{yang2025qwen3}, OLMo-2 7B and 13B \citep{olmo2024olmo2}, phi-4 \citep{abdin2024phi4}, seven Pythia models from 160 million to 12 billion parameters trained on the same data in the same order \citep{biderman2023pythia}, and six checkpoints of Pythia-1.4b from initialization to the end of training. Step 0 serves as the random-weight control. To ask whether recovered roles follow meaning, we pair each relation with a converse paraphrase that reverses the order of mention and keeps the meaning, such as ``a tulip on a barrel'' and ``a barrel under a tulip'' (Figure~\ref{fig:protocol}d). The tulip is the \emph{figure} in both. A readout of the figure, fitted on canonical sentences and scored on converse ones, has a transfer of 1 if the state codes semantic roles and 0 if it codes the order of mention.

\paragraph{Image encoders.} Parts are Tiny-ImageNet photographs \citep{le2015tiny}. We split the 200 classes into 120 for fitting, 30 for validation and 50 fresh, so fresh parts come from classes the readout never saw. Two photographs share a $224 \times 224$ canvas \textsc{beside} or \textsc{above} each other, or \textsc{overlap}ping with the second covering a corner of the first. We test CLIP ViT-B/32 and ViT-H/14 \citep{radford2021clip,cherti2023}, DINOv2 small, large and giant \citep{oquab2024dinov2}, and SigLIP so400m \citep{zhai2023siglip}, each with trained and random weights. A part's code is the encoder's state for the photograph alone. The pretrained models use 3 seeds, each a new random split of the parts.

%=====================================================================
\section{Results}
\label{sec:results}

In this section, we report seven results.
\begin{itemize}
  \item \textbf{R1.} Recovery of fresh parts separates role--filler composition from blends and superposition, and it meets the proven ceilings (Figure~\ref{fig:calibration}).
  \item \textbf{R2.} Networks trained on few glyphs keep fresh glyphs in their states and fail to render them (Figure~\ref{fig:networks}).
  \item \textbf{R3.} Trained networks are nearly role-factorized, and their decoders use the role codes (Figure~\ref{fig:structure}).
  \item \textbf{R4.} A learned composer stays invertible beyond its training depths, while its decomposer fails one level deeper (Figure~\ref{fig:depth}a,b).
  \item \textbf{R5.} Training lifts pretrained text encoders above a blend by very different amounts, and none approaches the role--filler reference (Figure~\ref{fig:depth}c,d).
  \item \textbf{R6.} Scale leaves language models and image encoders at most 60\% of the way to the reference, and recovery in Pythia-1.4b peaks early in training (Figure~\ref{fig:real}a,b,d).
  \item \textbf{R7.} Language models encode which noun is the figure, yet they bind roles to word order (Figure~\ref{fig:real}c).
\end{itemize}

\subsection{Fresh parts separate composition from blends}
\label{sec:calibration}

Given the candidates, a lookup identifies the role-1 part of every operator that is injective on the vocabulary with accuracy 1.00, including half-rank projections and HRR, and it fails only for the symmetric blends (0.52; Figure~\ref{fig:calibration}a). Recovery of fresh parts separates the operators and meets the ceiling of Proposition~\ref{prop:ceiling} (Figure~\ref{fig:calibration}b). At $V = 128$, linear readouts recover the role--filler operators and full-rank projections with $R^2 = 1.00$, half-rank projections with 0.47, $x + 0.5y$ with 0.80 against a ceiling of 0.80, and $x + y$ with 0.50. HRR and MAP reach 0.44 and 0.47, close to the same ceiling of $1/2$. A lookup thus decodes vector-symbolic superposition of dense parts over a vocabulary, while no readout recovers fresh parts from it. The readout itself needs enough parts (Figure~\ref{fig:calibration}c). With 8 parts in $\R^{16}$, least squares recovers fresh concatenated parts with $R^2 = 0.38$, and it reaches 1.00 from 24 parts, within 0.01 of the closed form of Proposition~\ref{prop:span}. MLP readouts transfer worse than linear ones, 0.11 against 0.47 at half rank, because a flexible readout memorizes the vocabulary (Figure~\ref{fig:calibration}d).

\begin{figure}[tbp]
  \centering
  \includegraphics[width=\linewidth]{fig_calibration.pdf}
  \caption{\textbf{Recovery of fresh parts separates role--filler composition from blends and superposition.} Fourteen operators combine Gaussian parts in $\R^{16}$ in pairs. \textbf{a.} Given the candidates, a lookup over their composites identifies the role-1 part of every operator that is injective on the vocabulary (blue, role--filler; grey, fixed-width superposition) and fails only for symmetric blends (orange); $V = 32$ training parts. \textbf{b.} A linear readout scored on fresh parts recovers the share that the ceiling of Proposition~\ref{prop:ceiling} allows; $V = 128$. \textbf{c.} Recovery of familiar (dashed) and fresh (solid) parts against the number of training parts. The dotted line gives the closed form of Proposition~\ref{prop:span} for least squares on concatenation, and we clip values below $-0.2$. \textbf{d.} MLP readouts transfer worse than linear ones, and we draw negative values at zero. Means over 5 seeds, with standard errors.}
  \label{fig:calibration}
\end{figure}

\subsection{Trained networks keep parts they cannot read}
\label{sec:networks}

Every architecture reproduces unseen pairs of training glyphs with output $R^2$ between 0.85 and 1.00, except the transformer trained on eight glyphs (Figure~\ref{fig:networks}b, dashed). Fresh glyphs separate the networks. At $m = 64$, their output $R^2$ rises from 0.20--0.34 at $V = 8$ to 0.85--0.91 at $V = 64$ (Figure~\ref{fig:networks}b, solid). The states tell a different story (Figure~\ref{fig:networks}c). A population readout recovers fresh glyphs from the MLP, CNN and role--filler states with $R^2$ of 0.83--0.86 at $V = 8$ and 0.89 at $V = 64$, against 0.82--0.83 for the untrained networks and 0.91 for the input pixels. Against concatenated (0.98) and summed (0.48) references, these values give $\kappa = 0.73$ at $V = 8$ and $0.83$ at $V = 64$. The glyph code is informative, because an affine map recovers the latent with $R^2 = 0.92$, and the residual $\delta$ is 0.02--0.04 (Figure~\ref{fig:structure}). At $V = 64$ the states are therefore representationally compositional at a tolerance of about $\tau = 0.17$.

A small vocabulary breaks the networks' use of their own states. A vocabulary readout recovers only 0.23--0.29 at $V = 8$, and it tracks each network's own fresh-glyph behavior better than the population readout, with $r = 0.50$ against $0.30$ within conditions and a 95\% bootstrap interval of 0.09 to 0.32 for the difference (Figure~\ref{fig:networks}d). A readout of the raw inputs adds nothing once the trained-state readout is known (Appendix~\ref{app:protocol}). The decoders thus behave like readouts fitted on their vocabulary, as Proposition~\ref{prop:span} predicts, and these networks fill the upper right cell of Table~\ref{tab:twobytwo}. The transformer is the exception, with a less recoverable state at $V = 8$ (0.59; Figure~\ref{fig:networks}c). Distinct parts matter beyond combinations. With 140 training pairs drawn from 16 to 64 glyphs, fresh-glyph behavior rises from 0.47--0.64 to 0.81--0.82, while behavior on held-out pairs falls by 0.05--0.10 (Figure~\ref{fig:networks}e). Width matters only below $kq = 16$, where the population readout reaches 0.67--0.71 at $m = 8$ against the Eckart--Young ceiling of 0.79 (Proposition~\ref{prop:ey}). The non-additive \textsc{occlude} relation lowers the population readout of every architecture from 0.85--0.89 to 0.75--0.83 (Figure~\ref{fig:networks}f).

\begin{figure}[tbp]
  \centering
  \includegraphics[width=\linewidth]{fig_networks.pdf}
  \caption{\textbf{Networks trained on few glyphs keep fresh glyphs in their states and fail to render them.} An encoder writes a state from a composite of two glyphs, and a decoder renders them under a target relation. \textbf{a.} One pair of glyphs under the three relations. \textbf{b.} Output $R^2$ on unseen pairs of training glyphs (dashed) and on fresh glyphs (solid) against the number of training glyphs $V$ (width $m = 64$, \textsc{over}). \textbf{c.} Recovery of fresh glyphs from the state by a population readout fitted on composites of 400 other glyphs (solid) and by a vocabulary readout fitted on the training composites (dashed). Dotted lines show untrained networks, and the black dashed line shows the input pixels. \textbf{d.} Each point is one run. The network's own fresh-glyph behavior follows the vocabulary readout (colored) more closely than the population readout (grey). \textbf{e.} With 140 training pairs, more distinct glyphs raise fresh-glyph behavior (solid) at a small cost on held-out pairs (dashed). \textbf{f.} Population readout against width for \textsc{over} (solid) and \textsc{occlude} (dashed), with the Eckart--Young ceiling ($V = 64$). Means over 9--10 seeds, with standard errors.}
  \label{fig:networks}
\end{figure}

\subsection{Decoders use the role codes}
\label{sec:structure}

We next ask whether the states are role-factorized and whether the decoders use the role codes. Under \textsc{over}, the residual $\delta$ is 0.02--0.04 for the MLP, CNN and role--filler networks and 0.06--0.13 for the transformer. Role maps predict the reversed state from the unreversed one with $R^2$ of 0.96--0.99, and the decoder renders the predicted reversals as well as the true ones, 0.86 against 0.86 at $V = 64$ (Figure~\ref{fig:structure}a). A training glyph placed in a role it never occupied, through its code read from another composite, reproduces the network's own output, 0.81 against 0.81 (Figure~\ref{fig:structure}b). The interchange intervention renders the donor glyph with $R^2 = 0.78$, against 0.81 when the network sees the donor itself, and it spills 0.16 into the untouched half of the output (Figure~\ref{fig:structure}c). The network's own response to a change of one glyph spills 0.25. Edits off the code distribution spill far more, 0.90 for the natural direction $h(a', b') - h(a, b)$, 0.95 for a random direction and 0.75 for a random direction in the range of $B_1$, against 0.30 when its coefficients follow the code distribution. The decoders therefore use the role components the readouts find, which answers the objection of \citet{fodor1990} for these networks. Under \textsc{occlude}, $\delta$ rises to 0.11--0.22, predicted reversals decode worse than true ones (0.56 against 0.70), and interventions spill 0.31.

Learning a new glyph exposes a limit of this structure (Figure~\ref{fig:structure}d). Fine-tuned on a fresh glyph in role~1 alone, networks trained on 64 glyphs render it in role~2 worse than before, 0.62 against 0.75, and their held-out pairs fall from 0.98 to 0.90. Replay prevents both losses. Writing the glyph's code into role~2 through $B_2$, with no gradient step, renders it at 0.75 with no cost to old pairs. A role-factorized state lets a system add a part by adding its code, and fine-tuning on one role ignores that structure.

\begin{figure}[tbp]
  \centering
  \includegraphics[width=\linewidth]{fig_structure.pdf}
  \caption{\textbf{The decoders of trained networks use the role codes that readouts find.} Role maps $h \approx c + B_1 e(x) + B_2 e(y)$, fitted on a pool of glyphs, predict states that each network's own decoder renders. \textbf{a.} Output $R^2$ for the role reversal of a fresh composite, decoded from the true reversed state, from a state predicted from the unreversed composite alone, and from the unreversed state ($V = 64$; we clip bars at $-0.5$). \textbf{b.} A training glyph in a role it never occupied, rendered by the network and from the predicted state. \textbf{c.} The spill of an edit, the change in the untouched half of the output relative to the edited half, for interventions on each role and for edits of the same norm. \textbf{d.} Old pairs and a fresh glyph in each role before and after fine-tuning on the glyph in one role, with replay, and in both roles, and after writing its code into role 2 (\textsc{over}, $V = 64$). Means over 5 seeds, four architectures and two widths, with standard errors.}
  \label{fig:structure}
\end{figure}

\subsection{Recursion breaks the decoder first}
\label{sec:depth}

Within its training depths, the composer's own decomposer recovers fresh leaves almost perfectly for $w \ge 32$ (Figure~\ref{fig:depth}a). At depth 3 it fails even at $w = 256$, with $R^2 = 0.26$, while a population readout recovers the same leaves with 0.99. At depths 4 and 5, which the composer never saw, the population readout stays within 0.14 of the ceiling of Proposition~\ref{prop:depth}. The relation stays invertible on its own outputs, and the decomposer fails to follow it. HRR shows the mirror case (Figure~\ref{fig:depth}b). A population readout recovers its leaves exactly whenever $w \ge 2^n d$, while recursive unbinding loses about $(2^n - 1)d/w$ of the variance to crosstalk. At $w = 256$ unbinding falls from 0.96 at depth 1 to $-0.13$ at depth 5, against a prediction of 0.97 and 0.03.

\subsection{Training lifts text encoders partway}
\label{sec:pretrained}

The pipeline recovers the concatenation reference with $R^2 = 1.00$ and the sum reference with 0.48 (Figure~\ref{fig:depth}c). On fresh nouns, the three GPT-2 models come closest to the concatenation reference with $\kappa = 0.44$--$0.50$, followed by CLIP-H with 0.36, and the other encoders reach 0.06--0.26. A role--filler reference that mixes the roles through a random full-rank map raises $\kappa$ by at most 0.06. The interaction share on a $12 \times 12$ grid of fitting nouns, a lower bound on $\delta$, is 0.007--0.012 for GPT-2 and at most 0.02 for CLIP. The same architectures with random weights reach $\kappa = -0.11$ to $0.15$, so training raises $\kappa$ by 0.51--0.61 for GPT-2, by 0.26--0.30 for T5 and CLIP-H, and by at most 0.17 for BERT and the two smaller CLIP encoders. Role assignment and identification among 100 candidates separate trained from random weights less well than recovery (Figure~\ref{fig:depth}d). Two caveats bound these values. The control ``a $x$ and a $y$'' gives the same recovery as the asymmetric relations, 0.75 against 0.74 for GPT-2 large, so the readouts recover the order of the nouns. The number of tokens in a noun dominates the codes of the two smaller CLIP encoders. Removing it drops them below their random-weight counterparts, 0.47 and 0.46 against 0.56 and 0.53, while GPT-2 and CLIP-H stay ahead with $\kappa = 0.35$--$0.42$. The encoder that conditions Stable Diffusion 1.5 thus gains little from training beyond the token count, consistent with the argument that contrastive training gives little incentive to encode word order \citep{yuksekgonul2023}.

\begin{figure}[tbp]
  \centering
  \includegraphics[width=\linewidth]{fig_depth_pretrained.pdf}
  \caption{\textbf{A recursive composer outlives its decomposer, and training lifts text encoders only partway.} \textbf{a.} Recovery of fresh leaves from a learned recursive composer of width $w$, trained on depths 1 and 2 (shaded), by its own decomposer (solid) and by a population readout (dashed), with the ceiling of Proposition~\ref{prop:depth} (dotted; $V = 16$; we clip values at $-0.5$). \textbf{b.} The same for holographic reduced representations, with the crosstalk prediction $1 - (2^n - 1)d/w$ (dotted). \textbf{c.} Recovery of nouns from sentence states on new pairs of fitting nouns (grey) and on fresh nouns (orange), against the concatenation (blue) and sum (dark grey) references built from each model's own part codes. Crosses and dots show random weights scored on their own and on the trained model's codes. \textbf{d.} Role assignment and identification among 100 candidates on fresh nouns. Dark grey marks show the sum reference, and crosses show random weights. Means over 5 seeds (a, b) and 3 seeds (c, d).}
  \label{fig:depth}
\end{figure}

\subsection{Scale leaves large models short of the reference}
\label{sec:llm}

We next ask whether larger and longer-trained models come closer. No language model approaches the concatenation reference (Figure~\ref{fig:real}a). Within Pythia, $\kappa$ rises from 0.27 at 160 million parameters to 0.44 at 1 billion and levels off at 0.56--0.58 for 6.9 and 12 billion. Qwen3-8B, OLMo-2 7B and 13B and phi-4 reach 0.38--0.47, no more than the GPT-2 models, so size orders models within one training recipe and fails to order them across recipes. The training trajectory of Pythia-1.4b rises and then falls (Figure~\ref{fig:real}b). Its $\kappa$ is $-0.02$ at initialization, peaks at 0.70 after 2{,}000 steps, 1.4\% of training, and declines to 0.57, 0.52 and 0.47 at steps 8{,}000, 32{,}000 and 143{,}000. The interaction share stays at or below 0.03 at every checkpoint, so the states stay additive throughout the decline.

\label{sec:vision}Image encoders show the same shortfall (Figure~\ref{fig:real}d). Every trained encoder beats the sum reference, with $\kappa$ from 0.17 for DINOv2-small to 0.41 for DINOv2-giant, 0.44 for CLIP-H and 0.49 for SigLIP. The trained encoders identify fresh photographs among 300 with accuracy 0.66--0.93, against 0.40--0.42 for the sum reference, and the same architectures with random weights fall to $\kappa = -0.29$ to $0.03$. Within CLIP and within DINOv2 the larger model recovers more, and the states stay nearly additive, with an interaction share of at most 0.08. The \textsc{overlap} relation gives the highest $\kappa$, 0.30--0.60, because its larger photograph is easy to recover (0.82--0.90).

\subsection{Language models bind roles to word order}
\label{sec:roles}

Every trained model except BERT-large assigns the first-mentioned noun to its role with accuracy 0.94--1.00, and BERT-large reaches 0.82 (Figure~\ref{fig:real}c). A readout of the figure fitted on both phrasings assigns it with 0.79--0.85 in the language models, against 0.54 for Pythia-1.4b at initialization, so the trained states carry which noun is the figure. Fitted on canonical sentences, the same readout assigns the figure of converse sentences with accuracy 0.00--0.05 in every language model, OLMo-2 13B and phi-4 included. The bidirectional T5-large and BERT-large transfer more, 0.20 and 0.32, and no model approaches 1. These models therefore code the order of mention in a role-factorized way. The semantic figure enters their states through codes that change with the phrasing, which violates role-factorization for semantic roles (Definition~\ref{def:recoverable}).

\begin{figure}[tbp]
  \centering
  \includegraphics[width=\linewidth]{fig_real.pdf}
  \caption{\textbf{Scale and training leave large models short of a role--filler reference, and language models bind roles to word order.} The score $\kappa$ is 1 for the concatenation of the part codes and 0 for their sum. \textbf{a.} $\kappa$ on fresh nouns against the number of parameters for the Pythia suite (line) and four other language models. \textbf{b.} $\kappa$ across the training of Pythia-1.4b, where step 0 is the initialization. \textbf{c.} Role assignment on fresh nouns by readouts of the first-mentioned noun (blue) and of the figure (orange), each fitted on both phrasings, and of the figure fitted on canonical sentences and scored on converse ones (red). A readout of word order scores 0 on the last test and a readout of meaning scores 1, and ``Py init'' is Pythia-1.4b at step 0. \textbf{d.} $\kappa$ on fresh photographs for image encoders with trained (blue) and random (grey) weights, averaged over three relations. Means over 3 seeds, with standard errors.}
  \label{fig:real}
\end{figure}

%=====================================================================
\section{Conclusion and limitations}
\label{sec:discussion}

We studied what makes a representation compositional by requiring that a fixed, declared readout recover every part in its role from composites of parts the readout never saw. Analytic operators met the proven ceilings, and recovery of fresh parts separated role--filler composition from blends and superposition. Small networks kept fresh parts in nearly role-factorized states, and their decoders used the role codes, yet training on few parts left decoders that could not read new parts. A learned composer stayed invertible beyond the depth at which its decomposer failed. In 28 pretrained models, training lifted every state above a blend, and no fully trained model got past 60\% of the way to a role--filler reference. Recovery in Pythia-1.4b peaked after 1.4\% of training, and every language model bound roles to word order. These findings suggest that a system's states and its procedures for reading them become compositional separately, and that scale and training leave both incomplete in the models we tested.

Compositionality matters because a finite experience must support an unbounded number of new situations. A behavioral test on unseen pairs of familiar parts cannot establish this ability, because a lookup over the vocabulary passes it. Our results suggest five requirements for a network. Its state must have room for every part (Propositions~\ref{prop:ceiling} and~\ref{prop:ey}). Its procedure for reading parts must learn from many distinct parts, because combinations of few parts leave the procedure undetermined (Proposition~\ref{prop:span}). Each part should enter every role through one code, so that one composite fixes its reversal (Proposition~\ref{prop:systematic}), and the language models fail this requirement for semantic roles. The role components must be the ones the system uses. The relation must stay invertible on its own outputs, and the system's procedure must follow it (Proposition~\ref{prop:depth}). We therefore argue that a claim that a model represents compositionally should state its readout class, its part codes and its fresh parts, and should report the model's own procedures separately.

\paragraph{Limitations.} Our trained networks are small and their domains rendered, chosen so that we could draw fresh parts. In pretrained models the parts are fresh only to the readout, and a richer class of methods, such as discovered tensor-product decompositions \citep{soulos2020}, might find a decomposition our declared readouts miss. Our tests of pretrained models measure their states, and the half of the definition that concerns a system's own procedures remains untested there. Human studies have rarely used parts fresh in our sense (Appendix~\ref{app:cases}), and a decoder of agent and patient \citep{frankland2015} fitted on some nouns and tested on others would give the direct neural analogue.
